\chapteropen{09}

\begin{cele}
\item policzyć roczny koszt czytania na trzy sposoby --- kupując, pożyczając i wymieniając --- i ułożyć zestaw, który mieści się w budżecie;
\item założyć klub czterech osób z zasadami, które przetrwają pierwsze miesiące, i zaplanować jego rok;
\item odróżnić wspólne czytanie, które podtrzymuje nawyk, od rywalizacji o liczbę książek, która go psuje.
\end{cele}

Czytanie bywa opisywane jako zajęcie samotne: fotel, lampa, ty i tekst. Samotność ma swoje zalety, ale ma też koszt, który rzadko się liczy. Nikt nie zapyta, dlaczego przerwałeś w połowie, i nie ma komu być winnym skończenia. Ten rozdział pokazuje, jak czytanie z innymi zmienia rachunek czasu i pieniędzy, i w którym miejscu wspólne czytanie zaczyna szkodzić.

\section{Biblioteka: zero złotych, ale nie zero czasu}

Biblioteka jest najtańszym sposobem czytania, jaki istnieje. Karta biblioteczna jest bezpłatna \verify{w większości bibliotek publicznych karta czytelnika jest bezpłatna}, a samo wypożyczenie nie kosztuje nic poza ewentualną opłatą za przetrzymanie. Koszty biblioteki są jednak inne niż koszty księgarni i łatwo je przeoczyć.

Pierwszy koszt to kolejka. Popularny tytuł, najnowsza powieść czy laureat ważnej nagrody, może czekać na swoją kolej tygodniami, bo tyle osób zamówiło go przed tobą. Nie ma tu nic złego, ale trzeba to uwzględnić w planie. Drugi koszt to termin zwrotu. Okres wypożyczenia jest ograniczony i zwykle wynosi kilka tygodni \verify{okres wypożyczenia w bibliotekach publicznych jest ograniczony i wynosi zwykle kilka tygodni}, a możliwość przedłużenia zależy od tego, czy ktoś inny nie czeka na tę samą pozycję. Książka odebrana w czwartek z trzytygodniowym terminem musi zostać przeczytana w tym terminie, a to zmienia tygodniowy plan.

Jak z tym pracować? Zamawiaj z wyprzedzeniem, a nie w dniu, w którym chcesz zacząć. Zamawiaj kilka tytułów naraz, bo ten, który przyjdzie pierwszy, nie zawsze jest tym, który chcesz czytać najpierw. Zapisz w kalendarzu termin zwrotu i ustal, ile dni przed nim zaczniesz czytać. Jeśli biblioteka nie ma danego tytułu, zapytaj o wypożyczenie międzybiblioteczne. W wielu sieciach jest to możliwe \verify{wypożyczenie międzybiblioteczne jest dostępne w wielu sieciach bibliotecznych}, ale czas oczekiwania bywa dłuższy niż przy zakupie z dostawą następnego dnia.

Biblioteka to też coś więcej niż magazyn. Bibliotekarka lub bibliotekarz zna zwykle katalog lepiej niż wyszukiwarka i potrafi wskazać coś z półki, o czym nie wiedziałeś. Dla czytelnika, który chce czytać dużo za mało pieniędzy, biblioteka jest pierwszą kolumną budżetu, a nie planem awaryjnym.

\section{Klub czterech osób}

Klub książkowy to najprostszy sposób, żeby czytanie stało się zobowiązaniem, nie zamieniając go w obowiązek. Zanim go założysz, zdecyduj o trzech rzeczach: wielkości, rytmie i sposobie wybierania książek.

\textbf{Wielkość.} Proponuję cztery osoby i nie traktuj tej liczby jako świętej. Przy czterech osobach każda wybiera książkę co cztery miesiące, więc rotacja mieści się w roku. Przy ośmiu osobach jeden wybór przypada co osiem miesięcy, zebranie zamienia się w zbiórkę opinii, a część członków przychodzi bez przeczytania. Przy dwóch osobach nie ma rotacji i rozmowa szybko się wyczerpuje. Ten dobór to sąd praktyczny, nie wynik badań. Traktuj go jako propozycję do zmiany, jeśli wasz kalendarz wygląda inaczej.

\textbf{Rytm.} Spotkanie raz w miesiącu, o tej samej porze i w tym samym miejscu. Stała godzina działa lepiej niż ustalanie każdego terminu osobno, bo każde ustalanie to okazja do odwołania. Zebranie trwa 90 minut. To wystarcza na wrażenia z książki, dwa pytania do grupy i wybór kolejnej pozycji. Dłużej zaczyna być spotkaniem towarzyskim, które ma tylko książkę w tle.

\textbf{Wybór.} Każdy członek wybiera jedną książkę w cyklu, a reszta czyta ją razem z nim. Wybierający odpowiada za pierwsze pytanie na spotkaniu. Dzięki temu wybór nie jest tylko ulubionym tytułem, lecz czymś, co wybierający chce bronić.

Zasady, które przetrwają pierwsze miesiące:
\begin{itemize}
\item Nieczytanie nie wyklucza ze spotkania. Przychodzisz i mówisz, na której stronie jesteś i dlaczego. To też jest wkład do rozmowy.
\item Porzucenie książki jest dozwolone, ale wymaga jednego zdania uzasadnienia. Bez tej zasady klub uczy się, że przyznanie się do porzucenia kosztuje, i wszyscy zaczynają udawać, że czytają.
\item Spotkanie zaczyna się od pytania o to, co zostało w pamięci, a nie od oceny. Pytanie \glqq czy była dobra\grqq{} kończy rozmowę. Pytanie \glqq co zostało z tobą po trzech dniach\grqq{} ją zaczyna.
\item Lista książek jest spisana na cały rok z góry, żeby nikt nie zgadywał, co czyta w marcu.
\end{itemize}

\begin{przyklad}[Harmonogram roku]
Czterech członków, dwanaście spotkań, każdy wybiera trzy książki w roku. Kolejność jest zapisana w tabeli z góry: styczeń wybiera Ania, luty Bartek, marzec Celina, kwiecień Darek, i tak dalej aż do grudnia. Każda osoba wybiera więc trzy razy, a dwanaście książek przechodzi przez cały klub.

Ile to wymaga czasu? Dwanaście książek po przeciętnie \val{typical.words} słów to 960\,000 słów rocznie na jednego członka. Przy 250 słowach na minutę wychodzi \val{need.12} minuty dziennie. To mniej niż przerwa na kawę, ale codziennie. Pojedyncza książka to około \val{typical.hours} godziny czytania, a przy czterech tygodniach między spotkaniami około 1,3 godziny tygodniowo.
\end{przyklad}

\section{Czytanie w dwoje i z dziećmi}

Czytanie we dwoje działa inaczej niż w klubie. Nie ma rotacji ani wyboru, jest jedna książka i dwie osoby, które spotykają się w tym samym miejscu, zwykle przy łóżku albo przy stole. Jego siła polega na tym, że to, co ważne dla jednej osoby, zostaje od razu zauważone przez drugą. Słowo \glqq zobacz, to ci się spodoba\grqq{} potrafi zmienić wieczór.

Przy dzieciach czytanie na głos jest powszechne i łatwe do zaplanowania. Pytanie, jaki dokładnie wpływ ma na rozwój języka dziecka, wymaga źródeł, których tu nie podaję \verify{czytanie na głos wspiera rozwój języka dziecka --- wielkość efektu wymaga przeglądu badań}. Dla dorosłego praktyczna korzyść jest prostsza: to stały, zaplanowany czas, który nie wymaga niczego ponad jedną książkę na tydzień i wspólne siedzenie w jednym miejscu.

Trzy zasady dla pary albo rodziny:
\begin{itemize}
\item Ustal jedną książkę na dwa tygodnie, a nie na miesiąc. Wspólne czytanie jest wolniejsze niż samotne, bo każde zdanie idzie przez dwie osoby, a plan z dłuższym terminem szybko przestaje obowiązywać.
\item Czytajcie w tych samych porach dnia, a nie wtedy, kiedy akurat wypadnie. Stała pora to jedyna część planu, której nie trzeba pilnować codziennie.
\item Jeśli czytasz na głos, planuj rozdziały krótsze, niż wydaje się rozsądne. Tempo wyznacza wtedy głos, a nie oczy, i łatwo przeliczyć się przy pierwszym dłuższym rozdziale.
\end{itemize}

Wspólne czytanie ma jeszcze jedną zaletę, którą łatwo pominąć: zostawia ślad w rozmowie. Po kilku tygodniach wiecie, co każde z was zapamiętało, i to jest materiał do powrotu do książki, który nie wymaga żadnej dyscypliny.

\section{Zobowiązanie i publiczny cel}

Publiczne ogłoszenie celu to popularna rada: napisz na blogu, powiedz znajomym, ustaw się w kolejce do rozliczenia. Ma sens, bo wiąże cię z opinią innych, ale ma też wadę. Deklaracja daje satysfakcję, zanim pojawi się jakakolwiek praca, i ta satysfakcja bywa tym, co zostaje. Nie znam rzetelnych danych, które by rozstrzygały, czy publiczna deklaracja zwiększa liczbę przeczytanych książek. Traktuj ją jako hipotezę do sprawdzenia we własnym dzienniku czytania, a nie jako gwarancję.

Zasada, która pomaga bardziej niż samo ogłoszenie: cel ma być sprawdzalny. \glqq Przeczytam dwadzieścia cztery książki\grqq{} da się sprawdzić w dzienniku w grudniu. \glqq Będę więcej czytać\grqq{} nie da się sprawdzić i dlatego nikt tego nie pamięta po trzech tygodniach. Jeśli chcesz zewnętrznego zobowiązania, wybierz osobę, która pyta o stan książki co tydzień, a nie tę, która słyszy o celu raz w styczniu. Zobowiązanie działa, dopóki jest tanie w utrzymaniu, czyli dopóki kosztuje jedno pytanie w tygodniu, a nie kolejne tłumaczenie się z wyniku.

\section{Budżet na książki}

Czytanie ma koszt roczny, który łatwo przeoczyć, bo rozkłada się na drobne zakupy. Poniższe liczby to przykład oparty na założeniach, które podaję wprost. Zamień je na swoje ceny.

Założenia: nowa książka kosztuje 50 zł, używana 15 zł, wypożyczenie z biblioteki 0 zł. Cel to 24 książki w roku, czyli dwie miesięcznie. Ceny nowych i używanych książek to moje przykładowe założenia, a nie dane rynkowe. Sprawdź je w swojej księgarni i na swojej giełdzie używanych książek, zanim cokolwiek zaplanujesz.

\begin{przyklad}[Cztery scenariusze na 24 książki]
Scenariusz A, same nowe: $24 \times 50$ zł = 1200 zł, czyli 1200 : 12 = 100 zł miesięcznie.

Scenariusz B, same używane: $24 \times 15$ zł = 360 zł, czyli 360 : 12 = 30 zł miesięcznie.

Scenariusz C, same z biblioteki: $24 \times 0$ zł = 0 zł. Koszt pieniężny znika, ale dochodzi kolejka i brak wyboru terminu.

Scenariusz D, mieszany: 4 nowe po 50 zł (200 zł), 12 używanych po 15 zł (180 zł) i 8 z biblioteki (0 zł). Razem 380 zł, czyli 380 : 24 = 15,83 zł na książkę i 380 : 12 = 31,67 zł miesięcznie.
\end{przyklad}

Scenariusz D zwykle się broni. Nowe kupujesz tylko te cztery, które chcesz mieć na półce, bo będziesz do nich wracać. Używane kupujesz te, które chcesz przeczytać, ale nie musisz zatrzymywać. Biblioteka przejmuje resztę, zwłaszcza tytuły, które są popularne i na które da się czekać.

Klub zmienia ten rachunek. Jeśli każdy z czterech członków kupuje swoje trzy wybory używane, to wydaje 3 $\times$ 15 zł = 45 zł rocznie, a pozostałe dziewięć tytułów czyta z cudzego egzemplarza, z biblioteki albo z wymiany. Każda osoba czyta wtedy dwanaście książek, a płaci za trzy.

Jedna reguła ogólna: jeśli miesięczny budżet na książki jest mniejszy niż cena jednej nowej książki, to nie jest to budżet, tylko loteria. Lepiej zaplanować 30 zł miesięcznie i wiedzieć, co za nie dostaniesz, niż wydawać nieregularnie i zastanawiać się, dlaczego nic nie zostało przeczytane.

\section{Wymiana książek i rozmowa jako nagroda}

Wymiana to najtańsza forma zakupu, jaką znam, ale ma swoje warunki. Na spotkaniu każdy przynosi jedną książkę, którą skończył i która mu nie zostaje na półce, i zabiera jedną. Zasada: nie wymieniasz książki, której nie przeczytałeś, bo wtedy nie masz nic do oddania i zostajesz bez wymiany. Druga zasada: oddajesz książkę w takim stanie, w jakim chciałbyś ją dostać. Ślady kawy i zagięte rogi to koszt, który płaci ktoś inny.

Wymiana przez pocztę ma koszt wysyłki, którego w przykładzie z budżetu nie uwzględniłem. Policz go, zanim zaczniesz. Dla książki o masie 400 gramów opłata za przesyłkę w obie strony może zjeść znaczną część oszczędności, więc wymiana na miejscu jest zwykle tańsza niż wymiana na odległość.

Nagrodą za skończoną książkę jest rozmowa o niej, a nie nowy przedmiot. Jeśli po każdym skończonym tytule kupujesz sobie następny, budżet z poprzedniej sekcji nigdy się nie domknie, a nagroda zamienia się w nawyk wydawania. Najlepszą nagrodą jest rozmowa, którą masz z kimś po przeczytaniu, i to ona najczęściej zostaje w pamięci dłużej niż przedmiot.

\section{Pułapka rywalizacji o liczbę}

Liczba przeczytanych książek jest mierzalna, łatwa do porównania i dlatego kusi. Klub, w którym ktoś prowadzi tabelę i publikuje ranking, zamienia czytanie w wynik. Znane od dawna prawo Goodharta mówi, że miara, która staje się celem, przestaje być dobrą miarą. Tutaj widać to bardzo wyraźnie.

\begin{falle}[Wygrana w tabeli]
Gdy liczy się liczbę książek, zaczynasz wybierać krótkie tytuły, bo dają punkty szybciej. Porzucasz trudną książkę, bo jej dwa miesiące nie wyglądają dobrze w rankingu, choć może była dla ciebie najważniejsza tego roku. W rozmowie mówisz tylko \glqq skończyłem\grqq{}, bo to jedyna część, która się liczy. Czytanie przestaje być rozmową o świecie, a zaczyna być rozmową o sobie.

Korekta jest prosta w zasadzie, choć trudna w praktyce. W klubie licz to, co zostało: rozmowy, porzucenia z uzasadnieniem i książki trudne, które wróciły do ciebie po miesiącu. Liczbę książek zapisuj w prywatnym dzienniku, gdzie możesz ją porównać z własnym poprzednim rokiem, a nie w tabeli, którą widzą inni. Wtedy liczba pozostaje informacją, a nie wyścigiem.
\end{falle}

Rywalizacja nie jest zła, jeśli wszyscy wiedzą, że to gra o regularność, a nie o wynik. Problem zaczyna się w chwili, gdy ktoś zaczyna czytać, żeby wygrać, a reszta zaczyna mu to wypominać.

\begin{cwiczenia}
\exitem{Ile kosztuje w roku 24 książki, jeśli 6 kupujesz nowe po 50 zł, a resztę wypożyczasz z biblioteki za darmo?}{$6 \times 50$ zł = 300 zł. Biblioteka nie dodaje kosztu, więc całość to 300 zł.}
\exitem{Klub czterech osób, każdy wybiera trzy książki w roku. Ile spotkań odbywa się w ciągu roku i ile książek przeczyta jeden członek?}{Spotkanie raz w miesiącu daje 12 spotkań. Każdy członek czyta wszystkie wybory, czyli 4 $\times$ 3 = 12 książek w roku.}
\exitem{Ile minut dziennie wymaga czytanie 12 książek po 80\,000 słów, jeśli czytasz 200 słów na minutę?}{12 $\times$ 80\,000 = 960\,000 słów. 960\,000 : 200 = 4800 minut w roku. 4800 : 365 $\approx$ 13,2, czyli około 13 minut dziennie.}
\exitem{Scenariusz D kosztuje 380 zł na 24 książki. Ile wynosi średnio jedna książka i ile miesięcznie?}{380 : 24 = 15,83 zł na książkę. 380 : 12 = 31,67 zł miesięcznie.}
\exitem{Wymień trzy zasady klubu, które dotyczą porzucania i nieczytania.}{Porzucenie jest dozwolone, ale wymaga jednego zdania uzasadnienia. Nieczytanie nie wyklucza ze spotkania, a przychodzący mówi, na której stronie jest i dlaczego. Spotkanie zaczyna się od pytania o zapamiętane, a nie od oceny.}
\exitem{Podaj jedną wartość, którą zapiszesz w dzienniku zamiast liczby książek w tabeli klubowej.}{Na przykład liczbę spotkań, na których powiedziałeś, co zostało z książki, albo liczbę porzuceń z uzasadnieniem. Obie mierzą rozmowę, nie wynik.}
\exitem{Osoba w klubie kupuje swoje trzy wybory używane po 15 zł. Ile wydaje rocznie i ile książek czyta?}{3 $\times$ 15 zł = 45 zł rocznie, a czyta 12 książek, bo pozostałe dziewięć pochodzi z cudzych egzemplarzy, biblioteki lub wymiany.}
\exitem{Co trzeba policzyć przed wymianą książek przez pocztę, czego nie ma w przykładzie z budżetu?}{Koszt wysyłki w obie strony dla każdej wymienianej książki, pomnożony przez liczbę wysyłek. Ten koszt nie jest w rachunku scenariuszy, więc dolicz go osobno.}
\end{cwiczenia}

\zadanie{W ciągu najbliższych 30 dni zaproś co najmniej jedną osobę do wspólnego czytania albo do klubu. Ustal termin pierwszego spotkania na 90 minut, zapisz trzy zasady z sekcji o klubie i spisz listę książek na najbliższe trzy miesiące. Miara sukcesu: w kalendarzu jest data kolejnego spotkania, a lista książek jest zapisana na papierze albo w notatniku, a nie tylko w głowie.}
